% ==============================================================================
% Section 3: Theoretical Framework
% ==============================================================================

\section{Theoretical Framework}
\label{sec:framework}

AR-DeC's design is grounded in two complementary theoretical frameworks: Keller's ARCS Motivational Model, which guides the gamification and engagement strategy, and Bloom's Revised Taxonomy, which structures the adaptive content delivery of the ITS engine. This section describes each framework and its mapping to specific system features.

\subsection{ARCS Motivational Model}
\label{subsec:arcs}

Keller's ARCS model \cite{keller2010} identifies four essential conditions for motivated learning: Attention, Relevance, Confidence, and Satisfaction. The model provides a systematic approach to instructional design by specifying strategies for establishing and maintaining each condition. AR-DeC maps each ARCS component to concrete system features, as summarised in \Cref{tab:arcs_mapping}.

\begin{table}[htbp]
\centering
\caption{Mapping of ARCS motivational model components to AR-DeC system features.}
\label{tab:arcs_mapping}
\begin{tabular}{p{2.2cm} p{4cm} p{6.5cm}}
\toprule
\textbf{ARCS Component} & \textbf{Strategy} & \textbf{AR-DeC Implementation} \\
\midrule
\textbf{Attention} &
Perceptual arousal; inquiry arousal; variability &
AR overlays with animated 3D models; novel visual annotations on scanned text; varied explanation formats (text, image, diagram); surprise badge unlocks \\
\addlinespace
\textbf{Relevance} &
Goal orientation; motive matching; familiarity &
Students scan their own textbooks and study materials; vocabulary is contextualised within the learner's actual reading; explanations reference domain-specific usage \\
\addlinespace
\textbf{Confidence} &
Learning requirements; success opportunities; personal control &
Adaptive difficulty via BKT-driven content selection; progress indicators (ProgressRing, mastery levels); learner controls explanation depth; quiz difficulty scales with demonstrated knowledge \\
\addlinespace
\textbf{Satisfaction} &
Intrinsic reinforcement; extrinsic rewards; equity &
Points for scanning (+5), reading (+3), quiz success (+10), daily streaks (+20); badge system with meaningful milestones; level progression (Beginner $\rightarrow$ Explorer $\rightarrow$ Scholar $\rightarrow$ Master); fair reward distribution tied to effort \\
\bottomrule
\end{tabular}
\end{table}

The ARCS model is particularly appropriate for AR-DeC because it bridges the gap between motivational psychology and practical instructional design. Unlike purely theoretical frameworks, ARCS provides actionable strategies that translate directly into system features and interaction patterns. The model's emphasis on maintaining motivation over time---rather than merely capturing initial interest---aligns with AR-DeC's goal of supporting sustained vocabulary acquisition across weeks of study.

\subsection{Bloom's Revised Taxonomy}
\label{subsec:blooms}

Bloom's original taxonomy of educational objectives \cite{bloom1956}, as revised by Anderson and Krathwohl \cite{anderson2001}, provides a hierarchical framework for classifying cognitive processes from lower-order (remembering, understanding) to higher-order (applying, analysing, evaluating, creating). AR-DeC's ITS engine uses this hierarchy to structure the adaptive delivery of vocabulary explanations and assessments, progressively challenging learners as their mastery increases.

The mapping between Bloom's taxonomy levels and AR-DeC's content delivery strategy is as follows:

\begin{enumerate}
    \item \textbf{Remember}: The system provides the basic definition of a scanned word, displayed as an AR text overlay. This is the default response for words the learner has not previously encountered.

    \item \textbf{Understand}: Upon subsequent encounters or when the student requests more detail, the system provides contextual examples, synonyms, and usage illustrations. AR overlays may include annotated diagrams or concept maps showing the word's relationship to related terms.

    \item \textbf{Apply}: The system presents the learner with fill-in-the-blank or sentence completion exercises using the target word in novel contexts, delivered through the quiz module.

    \item \textbf{Analyse}: For words with high mastery estimates, the system prompts the learner to identify relationships between the target word and previously learned vocabulary, compare usage across different contexts, or distinguish the word from semantically similar terms.
\end{enumerate}

The BKT student model determines which Bloom's level is appropriate for each word at each interaction point. Words with low mastery probability ($P(L_n) < 0.3$) receive Remember-level content; words with moderate mastery ($0.3 \leq P(L_n) < 0.6$) receive Understand-level explanations; words approaching mastery ($0.6 \leq P(L_n) < 0.85$) are assessed at the Apply level through quizzes; and words at high mastery ($P(L_n) \geq 0.85$) receive Analyse-level challenges to deepen and consolidate knowledge.

\subsection{Integration of Frameworks}
\label{subsec:framework_integration}

The ARCS model and Bloom's taxonomy operate synergistically within AR-DeC. The ARCS model governs the motivational and engagement layer---determining how content is presented and rewarded---while Bloom's taxonomy governs the cognitive layer---determining what content is delivered at each mastery level. Together, they ensure that learners receive appropriately challenging material (Confidence) presented in engaging formats (Attention) that are relevant to their current reading context (Relevance) and rewarded upon successful engagement (Satisfaction). This dual-framework approach addresses the criticism raised by Fowler \cite{fowler2015} that many educational technology applications lack sufficient theoretical grounding, and responds to the call by Ak\c{c}ay{\i}r and Ak\c{c}ay{\i}r \cite{akccayir2017} for stronger integration of learning theory with AR system design.

The theoretical framework is illustrated in \Cref{fig:framework}, which shows the mapping between ARCS components, Bloom's levels, and AR-DeC system features.

\begin{figure}[htbp]
    \centering
    % ==============================================================================
% TikZ Theoretical Framework Diagram
% ARCS Model + Bloom's Taxonomy mapped to AR-DeC Features
% ==============================================================================

\begin{tikzpicture}[
    node distance=0.6cm and 2cm,
    arcs/.style={
        rectangle,
        draw=#1!70!black,
        fill=#1!20,
        rounded corners=4pt,
        minimum width=3cm,
        minimum height=1cm,
        text centered,
        font=\small\bfseries,
        line width=1pt
    },
    feature/.style={
        rectangle,
        draw=gray!50,
        fill=gray!8,
        rounded corners=3pt,
        minimum width=4.5cm,
        minimum height=0.8cm,
        text centered,
        font=\scriptsize,
        line width=0.6pt
    },
    bloom/.style={
        rectangle,
        draw=violet!60!black,
        fill=violet!12,
        rounded corners=4pt,
        minimum width=2.5cm,
        minimum height=0.8cm,
        text centered,
        font=\small\bfseries,
        line width=1pt
    },
    connector/.style={
        -Stealth,
        thick,
        color=#1!60!black
    },
    heading/.style={
        font=\normalsize\bfseries\scshape,
        color=#1!70!black
    }
]

% --- Column Headers ---
\node[heading=blue] (h_arcs) at (0, 0) {ARCS Model};
\node[heading=gray] (h_features) at (5.5, 0) {AR-DeC Features};
\node[heading=violet] (h_bloom) at (11, 0) {Bloom's Level};

% --- ARCS Components (left column) ---
\node[arcs=red, below=1cm of h_arcs] (attention) {Attention};
\node[arcs=blue, below=0.8cm of attention] (relevance) {Relevance};
\node[arcs=green, below=0.8cm of relevance] (confidence) {Confidence};
\node[arcs=orange, below=0.8cm of confidence] (satisfaction) {Satisfaction};

% --- AR-DeC Features (center column) ---
\node[feature, below=1cm of h_features] (f1) {AR overlays \& 3D models};
\node[feature, below=0.3cm of f1] (f2) {Animated visual annotations};
\node[feature, below=0.6cm of f2] (f3) {Scan personal textbooks (OCR)};
\node[feature, below=0.3cm of f3] (f4) {Domain-contextual explanations};
\node[feature, below=0.6cm of f4] (f5) {Adaptive difficulty (BKT)};
\node[feature, below=0.3cm of f5] (f6) {Progress tracking \& mastery levels};
\node[feature, below=0.6cm of f6] (f7) {Points, badges, \& streaks};
\node[feature, below=0.3cm of f7] (f8) {Quiz combo multipliers};

% --- Bloom's Levels (right column) ---
\node[bloom, below=1cm of h_bloom] (remember) {Remember};
\node[bloom, below=0.8cm of remember] (understand) {Understand};
\node[bloom, below=0.8cm of understand] (apply) {Apply};
\node[bloom, below=0.8cm of apply] (analyse) {Analyse};

% --- ARCS to Features connections ---
\draw[connector=red] (attention.east) -- (f1.west);
\draw[connector=red] (attention.east) -- (f2.west);
\draw[connector=blue] (relevance.east) -- (f3.west);
\draw[connector=blue] (relevance.east) -- (f4.west);
\draw[connector=green] (confidence.east) -- (f5.west);
\draw[connector=green] (confidence.east) -- (f6.west);
\draw[connector=orange] (satisfaction.east) -- (f7.west);
\draw[connector=orange] (satisfaction.east) -- (f8.west);

% --- Features to Bloom's connections ---
\draw[connector=violet, densely dashed] (f1.east) -- (remember.west);
\draw[connector=violet, densely dashed] (f4.east) -- (understand.west);
\draw[connector=violet, densely dashed] (f5.east) -- (apply.west);
\draw[connector=violet, densely dashed] (f8.east) -- (analyse.west);

% --- Bloom's mastery thresholds ---
\node[font=\scriptsize\itshape, color=violet!60, right=0.2cm of remember] {$P(L_n) < 0.3$};
\node[font=\scriptsize\itshape, color=violet!60, right=0.2cm of understand] {$0.3 \leq P(L_n) < 0.6$};
\node[font=\scriptsize\itshape, color=violet!60, right=0.2cm of apply] {$0.6 \leq P(L_n) < 0.85$};
\node[font=\scriptsize\itshape, color=violet!60, right=0.2cm of analyse] {$P(L_n) \geq 0.85$};

\end{tikzpicture}

    \caption{Theoretical framework: integration of ARCS motivational model and Bloom's taxonomy in AR-DeC.}
    \label{fig:framework}
\end{figure}
